\section{Qunatum Efficiency Measurements}
\label{sec:Auswertung}

\subsection{Analysis}
Qunatum Efficiency (QE) of a Photovoltaic cell is defined as the ratio of the number of electron-hole pairs
created by the inceident light and the number of photons incident on the solar cell. The objective 
of the experiment was to estimate the spectral quantum efficiency, i.e., the wavelength dependant quantum
efficiency of different types of solar cells. 
\\

The measurement setup consisted of pyrodetector and samples. These are illuminated with light of a specific wavelength. With the use of monochromators and filters 
the wavelength is varied from $350~\si{\nano\meter}$ to $1150~\si{\nano\meter}$. To measure QE, two of the subsequent measurements are made:
\begin{itemize}
    \item A pyroelectric detector is used to determine the intensity of the light. This gives us the number of photons. 
    The pyrodetector delivers a voltage of $2878~\si{\volt}$ per $1~\si{\watt}$ of irradiated light.
    \item The pyroelectric detector is replaced by the solar cell and the photocurrent pro- vides information about the number of generated electron hole pairs.
\end{itemize}

To calculate QE from the measured values, a simple calculation to get the number of photons and the number of electron-hole pairs, needs to made. 
This calculation can be carried out by using, 
\begin{align*}
   N_{e-h} = \frac{I_{pc}t}{e} \\ 
   N_p = \frac{VAt \lambda}{2878 hc} \\ 
   QE = \frac{N_{e-h}}{N_p}
\end{align*}
where $N_{e-h}$ is the number of electron-hole pairs, $N_p$ is the number of photons, V is the measured voltage, $I_{pc}$ is the induced photocurrent
$A$ is the illuminated area, $e$ is the charge of the electron, $\lambda$ is the wavelength and $c$ is the speed of light. 
Fig \ref{fig:QE}, is the plot of calculated QE, for the various wavelengths. This by itself is not representative of the context of solar cells outside 
of the lab. To obtain this value, we normalise this plot to the value of $0.9$ for the wavelengths $300 < \lambda < 900$. 
Fig \ref{fig:NormalisedQE} represents the normalised curve, for the selected wavelengths. 
\begin{itemize}
\item The normalising factor for amorphous silicon is $27021.06$.
\item The normalising factor for crystalline silicon is $31978.62$.
\end{itemize}

As the final part of the analysis, we try to estimate the short circuit current expected for a solar spectrum for AM1.
This is done by using the following equation,
\begin{align*}
    I_{SC} = Aq\sum \frac{QE \cdot S(\lambda) \cdot \lambda \Delta\lambda}{hc}
\end{align*}

The short circuit current calculated by using the above equation for 

\begin{itemize}
\item Amorphous silicon is $1.88~\si{\milli\ampere}$
\item Crystalline silicon is $3.81~\si{\milli\ampere}$
\end{itemize}

\begin{figure}[H]
    \centering
    \includegraphics[width=\textwidth]{build/QE_Plots.pdf}
    \caption{Spectrum dependant Qunatum Efficiency}
    \label{fig:QE}
\end{figure}

\begin{figure}[H]
    \centering
    \includegraphics[width=\textwidth]{build/QE_NormalisePlots.pdf}
    \caption{Normalised Spectrum dependant Qunatum Efficiency}
    \label{fig:NormalisedQE}
\end{figure}


\subsection{Questions}
\textbf{1. Why does the quantum efficiency of the different cells drop to zero at 300 nm (different reasons for c-Si and a-Si solar cells)?}
Both a-Si and c-Si have different reason why the quantum efficiency drops to zero at $300\si{\nano\meter}$. 
For a-Si, the material has a short absorbtion depth for wavelengths in the specified range. This leads to all the photons being absorbed within the first few nanometers. However this regions sufffers from very 
poor electronic quality due to defects. This leads to high recombination losses that leads to the zero QE. 

For c-Si, on the other hand, the photons are all absorbed in the front layers of the cell, specifically near the anti-reflection coating, these layers are extremely close to the surface and not well-suited for efficient carrier collection.
So the created pairs suffer recombination losses before they reach the p-n junction. This prevents these pairs from contributing to the produced current. 

\textbf{2. Is there a correlation between the decrease of the quantum efficiency near the band edge and the the type of band gap of the absorber material?}
Yes there is a correlation between the decrease in the quantum efficiency near the band gap and the type of band gap.
In direct band gap materials, the absorption edge is sharp because photons can be absorbed directly without needing any momentum change; this means electron–hole pairs are created efficiently right at the band gap energy.
In indirect band gap materials, photon absorption near the band edge is much less efficient because it requires the involvement of a phonon to conserve momentum. This causes a much more gradual onset of absorption as photon energy approaches the band gap.
The result is a more gradual decrease in quantum efficiency near the cutoff